\scanpage{49}
% archive sheet number: 48

\mspage{12}
$\xi\in\underline D(\mathcal X)$, considérons (si $\varphi\colon\mathcal X\longrightarrow\underline\Psi$ \unsure{est} la projection)

\[
\begin{aligned}
\gamma_0&\overset{\deleted{\unsure{k}}\ \mathrm{déf}}{=}\varphi_*(\xi)\in\underline D(\underline\Psi)\\
\gamma&=\underline k_*(\gamma_0)\in D(\underline Q);
\end{aligned}
\]

donc

\[
\begin{aligned}
\gamma(0,0)&\simeq\deleted{\unsure{X_0}}H^\bullet_{\underline D}(X_0,\xi)\\
\gamma(0,1)&\simeq H^\bullet_{\underline D}(X_a,\xi)\\
\gamma(1,0)&\simeq H^\bullet_{\underline D}(X_b,\xi)\\
\gamma(1;1)&\simeq H^\bullet_{\underline D}(\mathcal X,\xi)
\end{aligned}
\]

et \deleted{diagramme} le carré squelettique associé est

$(\alpha)$\quad\diagram{\begin{tikzcd}[ampersand replacement=\&]
\& H^\bullet_{\underline D}(\mathcal X,\xi)\arrow[dl]\arrow[dr] \&\\
H^\bullet_{\underline D}(X_a,\xi)\arrow[dr] \& \& H^\bullet_{\underline D}(X_b,\xi)\arrow[dl]\\
\& H^\bullet_{\underline D}(X_0,\xi) \&
\end{tikzcd}.}

(NB\quad\diagram{\begin{tikzcd}[ampersand replacement=\&]
\& X_a\arrow[dr] \&\\
X_0\arrow[ur]\arrow[dr] \& \& \mathcal X\\
\& X_b\arrow[ur] \&
\end{tikzcd}}

n'a rien de commun, dans $\mathcal X$\ldots{} $X_a\cap X_b=\phi$. Mais

\diagram{\begin{tikzcd}[ampersand replacement=\&]
\& X_{/a}\arrow[dr,no head] \&\\
X_{/0}\arrow[ur,no head]\arrow[dr,no head] \& \& \mathcal X\\
\& X_{/b}\arrow[ur,no head] \&
\end{tikzcd}}

est comm., et le carré $(\alpha)$ est aussi déduit de
\mspage{13}
celui-là par $H^\bullet_{\underline D}(-,\xi)$.

L'élément $\gamma\in\underline D(Q)$ est cartésien par construction, donc le ``carré'' $(\alpha)$ \unsure{aussi|assoc.}. Ça serait le cas pour \uline{tt} cat. $\mathcal X$ \inserted{presque} au dessus de $\underline\Psi$, sans avoir à le supposer cofibré. Mais le fait que $\mathcal X$ soit cofibré doit donner au carré $(\alpha)$ la vertu suppl. qu'il est aussi \uline{cocartésien}. J'ai l'impression que ça ne peut absolument pas découler de l'un des axiomes Der 1 -- Der 5,5' vus jusqu'à présent. Je vais le mettre comme

\uuline{``Der 6''} \inserted{(provisoire)} \deleted{Soit} si $\mathcal X$ sur $\underline\Psi$ est cofibré, pour $\xi$ \deleted{const. et \unsure{dérivé|dérivée}, la} \inserted{quelconque} \deleted{\unsure{constant}} dans $\underline D(\mathcal X)$, \deleted{la} l'élément $\psi_*(\xi)\in\underline D(\underline Q)$, (où $\psi=(\mathcal X\xrightarrow{\varphi}\underline\Psi\lhook\joinrel\xrightarrow{\underline k}\underset{\substack{\shortparallel\\\Delta^1\times\Delta^1}}{\underline Q})$ est non seulement cartésien, mais aussi cocartésien.

% slanted
\margin{à vérifier que ça devient \unsure{déjà|clair} pour $\xi$ seulement (resp. const.), si \unsure{$D$ dérivé|$D$ cohomol.} \unsure{directe}.}

% slanted
\margin{\uline{Axiome des cofibrations} sous-entendu : si \unsure{le dérivateur|D de domaine} est orienté \unsure{en|comme} pour la codescente, où le ``bon'' axiome \unsure{(en th.?)|(\textnormal{en théorie})} est Der 6 de la page suiv.}
